\documentclass[11pt,letterpaper]{article}

\usepackage[margin=1in]{geometry}
\usepackage{amsmath,amssymb,amsthm,mathtools}
\usepackage{enumitem}
\usepackage{booktabs}
\usepackage{hyperref}
\usepackage{microtype}

\numberwithin{equation}{section}

\theoremstyle{plain}
\newtheorem{theorem}{Theorem}[section]
\newtheorem{lemma}[theorem]{Lemma}
\newtheorem{proposition}[theorem]{Proposition}
\newtheorem{corollary}[theorem]{Corollary}
\theoremstyle{definition}
\newtheorem{definition}[theorem]{Definition}
\newtheorem{example}[theorem]{Example}
\theoremstyle{remark}
\newtheorem{remark}[theorem]{Remark}
\newtheorem*{theoremA}{Theorem A}
\newtheorem*{theoremB}{Theorem B}
\newtheorem*{theoremC}{Theorem C}

\newcommand{\NN}{\mathbb{N}}
\newcommand{\ZZ}{\mathbb{Z}}
\DeclareMathOperator{\lcmop}{lcm}
\DeclareMathOperator{\Val}{Val}
\DeclareMathOperator{\im}{Im}

\title{First-appearance spectra and prime-power ladders for the Fibonacci modulus chain}
% Fill in author names, affiliations, corresponding-author e-mail addresses,
% and ORCID iDs before submission.
% Example:
% \author{First Author\\Department, University, City, Country\\
% \texttt{first.author@example.edu}\\ORCID: 0000-0000-0000-0000
% \and
% Second Author\\Department, University, City, Country\\
% \texttt{second.author@example.edu}\\ORCID: 0000-0000-0000-0000}
\author{}
\date{}

\begin{document}
\maketitle

\begin{abstract}
We study first-appearance phenomena for cyclic quotients along the Fibonacci
modulus chain $\ZZ/(F_{m+2}\ZZ)$ and its canonical realization on admissible
Zeckendorf words. We prove that a quotient onto $\ZZ/(q\ZZ)$ occurs at depth
$m$ if and only if $q\mid F_{m+2}$; equivalently, the least supporting depth is
$\alpha(q)-2$, where $\alpha(q)$ is the rank of apparition of $q$ in the
Fibonacci sequence. This gives the minimal common depth for any finite family
of moduli and an explicit formula for the fibers of the joint quotient map. We
next show that the numbers of cyclic, prime, and prime-power quotients born for
the first time at index $n$ are given by M\"obius transforms of $\tau(F_n)$,
$\omega(F_n)$, and $\Omega(F_n)$. Using Lengyel's valuation theorem, we compute
$\alpha(p^r)$ exactly and obtain rigid $p$-primary ladders. In particular, for
odd primes $p\neq 5$, the primitive $p$-layer already supports a quotient
modulo $p^2$ if and only if $p$ is a Wall--Sun--Sun prime. Finally,
reconstruction from proper quotient data exhibits a sharp dichotomy: levels
whose modulus has at least two distinct prime divisors are recovered from two
proper shadows, whereas prime-power levels require all prime-power factors of
$F_{m+2}$.
\end{abstract}

\medskip
\noindent\textbf{Keywords.} Fibonacci numbers, rank of apparition, M\"obius
inversion, Wall--Sun--Sun primes, Chinese remainder theorem.

\smallskip
\noindent\textbf{Mathematics Subject Classification (2020).}
11B39, 11A07, 11A25.

\section{Introduction and main results}

Let $\alpha(q)$ denote the rank of apparition of an integer $q\ge 2$ in the
Fibonacci sequence, that is, the least index $n$ for which $q\mid F_n$. The
starting point of the paper is the following first-appearance problem.
For which integers $q$ does the cyclic ring $\ZZ/(q\ZZ)$ occur naturally along
the Fibonacci modulus chain
\[
\ZZ/(F_3\ZZ),\ \ZZ/(F_4\ZZ),\ \ZZ/(F_5\ZZ),\ \dots,
\]
and at which index does it appear for the first time? We also ask for exact
counts of the quotients born at each index, for the structure of the
$p$-primary part of the hierarchy, and for the amount of quotient data needed
to reconstruct one level from proper residue information.

The chain above has a convenient model in Zeckendorf numeration. For each
$m\ge 1$, let
\[
X_m:=\{(w_1,\dots,w_m)\in\{0,1\}^m:\ w_iw_{i+1}=0\text{ for }1\le i<m\},
\]
the set of admissible Fibonacci words of length $m$. The finite value map
\[
V_m(w_1,\dots,w_m):=\sum_{k=1}^m w_kF_{k+1}
\]
identifies $X_m$ with $\{0,1,\dots,F_{m+2}-1\}$ by the finite-range form of
Zeckendorf's theorem. Pulling back addition and multiplication modulo
$F_{m+2}$ therefore equips $X_m$ with a finite commutative ring structure
canonically isomorphic to $\ZZ/(F_{m+2}\ZZ)$. The role of Zeckendorf words in
what follows is therefore structural rather than foundational: they provide a
uniform realization of the arithmetic questions on the Fibonacci modulus chain.

Classical apparition theory already supplies the basic divisibility input.
Wall and Vinson show that $q\mid F_n$ is controlled exactly by $\alpha(q)$
\cite{Wall1960FibonacciModm,Vinson1963RankApparition}. Primitive divisor
theorems explain when a genuinely new prime first appears
\cite{Yabuta2001PrimitiveDivisorFibonacci,BiluHanrotVoutier2001PrimitiveDivisorsLucas},
while Lengyel's valuation formula determines the $p$-adic depth of Fibonacci
divisibility \cite{Lengyel1995OrderFibonacciLucas}. Taken together, these
ingredients lead to a hierarchy theory with exact
first-appearance spectra for cyclic and prime-power quotients, a precise
$p$-primary ladder law, and a Wall--Sun--Sun criterion expressed at the
primitive layer itself. This is distinct from the algorithmic literature on
Zeckendorf addition and normalization
\cite{Frougny1988LinearNumerationOrderTwo,Frougny1992NumbersAutomata,Frougny1999OnlineAddition,
AhlbachUsatineFrougnyPippenger2013,LabbeLepsova2023FibonacciTwosComplement},
where the main questions concern implementation and normalization rather than
first occurrence and fixed-depth reconstruction.

\subsection*{Arithmetic form of the main theorems}

The core results may be read without the language of admissible words.

\begin{theoremA}
For $n\ge 2$, define
\[
A(n)=\#\{q\ge 2:\alpha(q)=n\},
\quad
A_{\mathrm{pr}}(n)=\#\{p\text{ prime}:\alpha(p)=n\},
\]
\[
A_{\mathrm{pp}}(n)=\#\{p^r:r\ge 1,\ p\text{ prime},\ \alpha(p^r)=n\}.
\]
Then
\[
A(n)=\sum_{d\mid n}\mu\!\left(\frac{n}{d}\right)\tau(F_d),\qquad
A_{\mathrm{pr}}(n)=\sum_{d\mid n}\mu\!\left(\frac{n}{d}\right)\omega(F_d),
\]
\[
A_{\mathrm{pp}}(n)=\sum_{d\mid n}\mu\!\left(\frac{n}{d}\right)\Omega(F_d).
\]
In particular, the exceptional indices $6$ and $12$ are prime-field empty but
not prime-power empty.
\end{theoremA}

\begin{theoremB}
For every prime $p$ and $r\ge 1$, the rank of apparition of $p^r$ is explicit.
One has $\alpha(5^r)=5^r$ and
\[
\alpha(2^r)=3\cdot 2^{\max(1,r-2)}\qquad (r\ge 2).
\]
If $p\neq 2,5$ and $e_p=\nu_p(F_{\alpha(p)})$, then
\[
\alpha(p^r)=\alpha(p)\,p^{\max(0,r-e_p)}.
\]
Hence, once $r\ge e_p$, the sequence $\alpha(p^r)$ grows by a factor of $p$ at
each step.
\end{theoremB}

\begin{theoremC}
Let $p\neq 2,5$ be prime. The following are equivalent:
\[
\alpha(p^2)=\alpha(p),\qquad \nu_p(F_{\alpha(p)})\ge 2,
\qquad p\text{ is a Wall--Sun--Sun prime}.
\]
Equivalently, the primitive $p$-layer is already locally thickened by a
quotient modulo $p^2$.
\end{theoremC}

Our first result is a complete classification of cyclic quotients at fixed
depth. Theorem~\ref{thm:quotient-observable-threshold} shows that there is a
unit-preserving quotient
\[
X_m\twoheadrightarrow \ZZ/(q\ZZ)
\]
if and only if $q\mid F_{m+2}$. If $\alpha(q)$ denotes the rank of apparition
of $q$, then the smallest supporting depth is $m_{\min}(q)=\alpha(q)-2$.
Theorems~\ref{thm:finite-signature-unified-resolution-minimal}
and~\ref{thm:finite-signature-joint-readout-image-fiber} extend this to finite
families of moduli.

The next step is to count the quotients that are born for the first time at a
given index. Theorem~\ref{thm:birth-spectrum-mobius} expresses the cyclic,
prime, and prime-power birth spectra as M\"obius transforms of $\tau(F_n)$,
$\omega(F_n)$, and $\Omega(F_n)$. In particular, the exceptional indices $6$
and $12$ are not empty layers: they carry no new prime fields, but they do
carry new prime-power quotients. Combined with the primitive divisor theorem,
this also shows that every resolution except $m=4$ and $m=10$ carries a prime
field quotient that first appears there; see
Theorem~\ref{thm:primitive-field-quotients}.

The local $p$-primary structure is described by
Theorem~\ref{thm:prime-power-apparition-ladders}. Using Lengyel's valuation
theorem, we determine $\alpha(p^r)$ for every prime power. For odd primes
$p\neq 5$, if $e_p=\nu_p(F_{\alpha(p)})$, then the primitive $p$-birth layer
already contains
\[
\ZZ/(p\ZZ),\ \ZZ/(p^2\ZZ),\ \dots,\ \ZZ/(p^{e_p}\ZZ),
\]
and thereafter the birth depths scale geometrically by a factor of $p$.
Corollary~\ref{cor:wss-thickening-criterion} interprets the Wall--Sun--Sun
condition as the presence of a quotient modulo $p^2$ already on the primitive
$p$-layer.

The final part of the paper concerns simultaneous quotient maps and exact
reconstruction. Theorems~\ref{thm:proper-shadow-compressibility}
and~\ref{thm:prime-power-shadow-complexity} show that CRT-split levels are
recovered from two proper quotient shadows, whereas prime-power levels require
all prime-power factors of $F_{m+2}$.

The paper is organized as follows. Section~2 fixes the ring structure on
$X_m$ and recalls the Fibonacci divisibility facts used later. Section~3
contains the exact threshold theorem and the M\"obius formulas for the birth
spectra. Section~4 treats prime-power ladders and the Wall--Sun--Sun criterion.
Section~5 studies finite signatures, reconstruction, and shadow complexity.
Appendix~A records a short table of initial birth data and ladder values.

\section{Preliminaries on Fibonacci divisibility and finite quotient rings}

For each $m\ge 1$, we continue to write
\[
X_m=\{(w_1,\dots,w_m)\in\{0,1\}^m:\ w_iw_{i+1}=0\text{ for }1\le i<m\}
\]
and
\[
V_m(w_1,\dots,w_m)=\sum_{k=1}^m w_kF_{k+1}.
\]

\begin{lemma}
\label{lem:zeckendorf-range-bijection}
For every $m\ge 1$, the map $V_m$ is a bijection from $X_m$ onto
$\{0,1,\dots,F_{m+2}-1\}$.
\end{lemma}

\begin{proof}
This is the finite-range form of Zeckendorf's theorem \cite{Zeckendorf1972}.
Every integer $N<F_{m+2}$ has a unique Zeckendorf expansion whose largest
nonzero digit occurs at a position at most $m$, and this expansion is exactly a
word in $X_m$.
\end{proof}

\begin{definition}
\label{def:finite-resolution-arithmetic}
For $w,v\in X_m$, define
\[
w\boxplus_m v:=w_m\bigl((V_m(w)+V_m(v))\bmod F_{m+2}\bigr),
\]
\[
w\boxtimes_m v:=w_m\bigl((V_m(w)V_m(v))\bmod F_{m+2}\bigr),
\]
where $w_m(N)$ denotes the unique word in $X_m$ with value
$N\in\{0,1,\dots,F_{m+2}-1\}$.
\end{definition}

\begin{proposition}
\label{thm:finite-resolution-mod}
For every $m\ge 1$, the value map gives a ring isomorphism
\[
V_m:(X_m,\boxplus_m,\boxtimes_m)\longrightarrow \ZZ/(F_{m+2}\ZZ).
\]
\end{proposition}

\begin{proof}
By Lemma~\ref{lem:zeckendorf-range-bijection}, $V_m$ is a bijection between
$X_m$ and a complete residue system modulo $F_{m+2}$. By construction,
\[
V_m(w\boxplus_m v)\equiv V_m(w)+V_m(v)\pmod{F_{m+2}}
\]
and
\[
V_m(w\boxtimes_m v)\equiv V_m(w)V_m(v)\pmod{F_{m+2}}.
\]
Thus $V_m$ is a bijective ring homomorphism.
\end{proof}

\begin{corollary}
\label{cor:phase-trichotomy}
Let $N:=F_{m+2}$.
\begin{enumerate}[label=(\arabic*)]
\item If $N$ is prime, then $X_m$ is a field.
\item If $N$ is a prime power of exponent at least $2$, then $X_m$ is a local
ring.
\item If $N$ has at least two distinct prime factors, then $X_m$ splits as a
nontrivial product of residue rings by the Chinese remainder theorem.
\end{enumerate}
\end{corollary}

\begin{proof}
This is immediate from Proposition~\ref{thm:finite-resolution-mod} and the
corresponding structure theory of $\ZZ/(N\ZZ)$.
\end{proof}

\begin{definition}
\label{def:rank-of-apparition}
For an integer $q\ge 2$, the \emph{rank of apparition} $\alpha(q)$ is the least
positive integer $n$ such that $q\mid F_n$.
\end{definition}

\begin{proposition}
\label{prop:rank-of-apparition-basic}
Let $q\ge 2$ and $n\ge 1$.
\begin{enumerate}[label=(\arabic*)]
\item One has $q\mid F_n$ if and only if $\alpha(q)\mid n$.
\item For any finite family $q_1,\dots,q_s\ge 2$,
\[
\alpha(\lcmop(q_1,\dots,q_s))=\lcmop(\alpha(q_1),\dots,\alpha(q_s)).
\]
\end{enumerate}
\end{proposition}

\begin{proof}
Part (1) is classical; see Wall \cite{Wall1960FibonacciModm} and Vinson
\cite{Vinson1963RankApparition}. For part (2), write $Q=\lcmop(q_1,\dots,q_s)$.
Then for every $n\ge 1$,
\[
Q\mid F_n
\iff
q_i\mid F_n\ \text{for all }i
\iff
\alpha(q_i)\mid n\ \text{for all }i
\iff
\lcmop(\alpha(q_1),\dots,\alpha(q_s))\mid n.
\]
Applying the minimality characterization in part (1) gives the formula.
\end{proof}

\section{Exact quotient thresholds and birth spectra}

\begin{theorem}
\label{thm:quotient-observable-threshold}
Let $m\ge 1$ and $q\ge 2$. The following are equivalent.
\begin{enumerate}[label=(\arabic*)]
\item There exists a unit-preserving quotient
\[
X_m\twoheadrightarrow \ZZ/(q\ZZ).
\]
\item $q\mid F_{m+2}$.
\item $\alpha(q)\mid (m+2)$.
\end{enumerate}
When these conditions hold, the quotient is unique. In particular,
\[
m_{\min}(q)=\alpha(q)-2.
\]
\end{theorem}

\begin{proof}
By Proposition~\ref{thm:finite-resolution-mod}, giving a unit-preserving
quotient on $X_m$ is equivalent to giving a unit-preserving surjective ring
homomorphism
\[
\widetilde{\pi}:\ZZ/(F_{m+2}\ZZ)\twoheadrightarrow \ZZ/(q\ZZ).
\]
Such a map, if it exists, is necessarily the reduction map
$\overline{t}\mapsto [t]_q$, and it is well defined and surjective exactly when
$q\mid F_{m+2}$. This proves (1) $\Leftrightarrow$ (2). The equivalence of
(2) and (3) is Proposition~\ref{prop:rank-of-apparition-basic}(1) with
$n=m+2$. The minimality statement is the case $n=\alpha(q)$.
\end{proof}

\begin{definition}
\label{def:birth-spectra}
For each $n\ge 2$, define
\[
A(n):=\#\{q\ge 2:\ \alpha(q)=n\},
\]
\[
A_{\mathrm{pr}}(n):=\#\{p\text{ prime}:\ \alpha(p)=n\},
\]
\[
A_{\mathrm{pp}}(n):=\#\{p^r:\ p\text{ prime},\ r\ge 1,\ \alpha(p^r)=n\}.
\]
\end{definition}

\begin{theorem}
\label{thm:birth-spectrum-mobius}
For every $n\ge 2$,
\[
A(n)=\sum_{d\mid n}\mu\!\left(\frac{n}{d}\right)\tau(F_d),
\]
\[
A_{\mathrm{pr}}(n)=\sum_{d\mid n}\mu\!\left(\frac{n}{d}\right)\omega(F_d),
\]
\[
A_{\mathrm{pp}}(n)=\sum_{d\mid n}\mu\!\left(\frac{n}{d}\right)\Omega(F_d).
\]
\end{theorem}

\begin{proof}
By Proposition~\ref{prop:rank-of-apparition-basic}(1), a nontrivial divisor
$q$ of $F_n$ satisfies $\alpha(q)\mid n$. Hence the nontrivial divisors of
$F_n$ are partitioned by their first-appearance index, and
\[
\tau(F_n)-1=\sum_{\substack{d\mid n\\ d\ge 2}} A(d).
\]
Since $n\ge 2$, one has
$\sum_{d\mid n}\mu(n/d)=0$, so the constant term $1$ disappears under M\"obius
inversion and
\[
A(n)=\sum_{d\mid n}\mu\!\left(\frac{n}{d}\right)\tau(F_d).
\]

The same argument with prime divisors instead of all divisors gives
\[
\omega(F_n)=\sum_{d\mid n} A_{\mathrm{pr}}(d),
\]
and with prime-power divisors gives
\[
\Omega(F_n)=\sum_{d\mid n} A_{\mathrm{pp}}(d).
\]
Applying M\"obius inversion to these two identities yields the remaining
formulas.
\end{proof}

\begin{corollary}
\label{cor:birth-spectrum-exceptional}
The exceptional indices $6$ and $12$ satisfy
\[
A_{\mathrm{pr}}(6)=A_{\mathrm{pr}}(12)=0,
\qquad
A_{\mathrm{pp}}(6)=A_{\mathrm{pp}}(12)=2.
\]
\end{corollary}

\begin{proof}
Since $F_6=8$, the only prime divisor is $2$, and $\alpha(2)=3$, so no new
prime field is born at $6$. On the other hand, $\alpha(4)=\alpha(8)=6$, hence
$A_{\mathrm{pp}}(6)=2$.

Likewise $F_{12}=144=2^4\cdot 3^2$. Its prime divisors are $2$ and $3$, with
$\alpha(2)=3$ and $\alpha(3)=4$, so $A_{\mathrm{pr}}(12)=0$. The new prime
powers are $9$ and $16$, because $\alpha(9)=12$ and $\alpha(16)=12$, whereas
$\alpha(4)=\alpha(8)=6$.
\end{proof}

\begin{theorem}
\label{thm:primitive-field-quotients}
For every $n\ge 3$ with $n\notin\{6,12\}$, there exists a prime $p_n$ such
that $\alpha(p_n)=n$. Equivalently, for every $m\ge 1$ with $m\notin\{4,10\}$,
the ring $X_m$ admits a quotient onto a prime field that does not factor
through any smaller depth.
\end{theorem}

\begin{proof}
For $n=3,4,5$ one may take $p_3=2$, $p_4=3$, and $p_5=5$. For $n\ge 7$ with
$n\neq 12$, the primitive divisor theorem for Fibonacci numbers
\cite{Yabuta2001PrimitiveDivisorFibonacci,BiluHanrotVoutier2001PrimitiveDivisorsLucas}
gives a prime $p_n$ dividing $F_n$ but no earlier Fibonacci number $F_d$ with
$1\le d<n$. By definition of rank of apparition, $\alpha(p_n)$ is the least
index at which $p_n$ appears. Since $\alpha(p_n)\mid n$ and $p_n$ divides no
earlier $F_d$, necessarily $\alpha(p_n)=n$.

The equivalent statement for resolutions follows from
Theorem~\ref{thm:quotient-observable-threshold}: $X_{n-2}$ admits a quotient
onto $\ZZ/(p_n\ZZ)$, and any factorization through $X_k$ with $k<n-2$ would
force $p_n\mid F_{k+2}$, hence $\alpha(p_n)\le k+2<n$, a contradiction.
\end{proof}

\section{Prime-power ladders and the Wall--Sun--Sun criterion}

\begin{theorem}
\label{thm:prime-power-apparition-ladders}
Let $p$ be prime and $r\ge 1$.
\begin{enumerate}[label=(\arabic*)]
\item For $p=5$,
\[
\alpha(5^r)=5^r.
\]
\item One has
\[
\alpha(2)=3,
\qquad
\alpha(2^r)=3\cdot 2^{\max(1,r-2)}\quad (r\ge 2).
\]
\item If $p\neq 2,5$ and $e_p:=\nu_p(F_{\alpha(p)})$, then
\[
\alpha(p^r)=\alpha(p)\,p^{\max(0,r-e_p)}.
\]
\end{enumerate}
\end{theorem}

\begin{proof}
For $p\neq 2,5$, Lengyel's valuation formula \cite{Lengyel1995OrderFibonacciLucas}
states that
\[
\nu_p(F_n)=
\begin{cases}
\nu_p(n)+\nu_p(F_{\alpha(p)}), & \alpha(p)\mid n,\\
0, & \alpha(p)\nmid n.
\end{cases}
\]
Since $\alpha(p)\mid p\pm 1$, one has $p\nmid \alpha(p)$. Writing
$n=\alpha(p)u$, the condition $p^r\mid F_n$ is equivalent to
\[
\nu_p(u)+e_p\ge r.
\]
The least such $u$ is $p^{\max(0,r-e_p)}$, which proves part (3).

The case $p=5$ follows from the explicit valuation formula
$\nu_5(F_n)=\nu_5(n)$, so the least $n$ with $5^r\mid F_n$ is $5^r$. For
$p=2$, the classical valuation formula gives the stated values
\cite{Lengyel1995OrderFibonacciLucas}.
\end{proof}

\begin{corollary}
\label{cor:prime-power-geometric-recursion}
Let $p\neq 2,5$ be prime and write $e_p=\nu_p(F_{\alpha(p)})$. If $r\ge e_p$,
then
\[
\alpha(p^{r+1})=p\,\alpha(p^r).
\]
Equivalently,
\[
m_{\min}(p^{r+1})+2=p\bigl(m_{\min}(p^r)+2\bigr).
\]
\end{corollary}

\begin{proof}
This is immediate from Theorem~\ref{thm:prime-power-apparition-ladders}(3).
\end{proof}

\begin{corollary}
\label{cor:wss-thickening-criterion}
Let $p\neq 2,5$ be prime. The following are equivalent.
\begin{enumerate}[label=(\roman*)]
\item $\alpha(p^2)=\alpha(p)$.
\item The primitive $p$-birth layer already carries a quotient
\[
X_{\alpha(p)-2}\twoheadrightarrow \ZZ/(p^2\ZZ).
\]
\item $\nu_p(F_{\alpha(p)})\ge 2$.
\item $p$ is a Wall--Sun--Sun prime.
\end{enumerate}
\end{corollary}

\begin{proof}
By Theorem~\ref{thm:prime-power-apparition-ladders}(3),
\[
\alpha(p^2)=\alpha(p)
\iff
\nu_p(F_{\alpha(p)})\ge 2.
\]
The equivalence of (i) and (ii) is exactly
Theorem~\ref{thm:quotient-observable-threshold}. For $p\neq 2,5$, the
Wall--Sun--Sun condition is equivalent to $\nu_p(F_{\alpha(p)})\ge 2$; see
Sun and Sun \cite{SunSun1992FibonacciFLT}.
\end{proof}

\section{Finite signatures, reconstruction, and shadow complexity}

\begin{definition}
\label{def:joint-readout}
Let $S=\{q_1,\dots,q_s\}$ be a finite family of integers at least $2$, and let
$m\ge 1$ satisfy $q_i\mid F_{m+2}$ for all $i$. The \emph{joint readout} at
depth $m$ is
\[
\rho_{m;S}:X_m\longrightarrow \prod_{i=1}^s \ZZ/(q_i\ZZ),
\qquad
x\longmapsto \bigl([V_m(x)]_{q_1},\dots,[V_m(x)]_{q_s}\bigr).
\]
\end{definition}

\begin{theorem}
\label{thm:finite-signature-unified-resolution-minimal}
Let $S=\{q_1,\dots,q_s\}$ and write $Q=\lcmop(q_1,\dots,q_s)$. The least depth at
which all moduli in $S$ occur simultaneously is
\[
L(S)=\alpha(Q)-2.
\]
\end{theorem}

\begin{proof}
By Theorem~\ref{thm:quotient-observable-threshold}, all $q_i$ occur at depth
$m$ if and only if $\alpha(q_i)\mid m+2$ for all $i$. By
Proposition~\ref{prop:rank-of-apparition-basic}(2), this is equivalent to
$\alpha(Q)\mid m+2$. The minimal such $m$ is therefore $\alpha(Q)-2$.
\end{proof}

\begin{theorem}
\label{thm:finite-signature-joint-readout-image-fiber}
Let $m\ge 1$, let $S=\{q_1,\dots,q_s\}$ with $q_i\mid F_{m+2}$ for all $i$, and
write
\[
M:=F_{m+2},
\qquad
Q:=\lcmop(q_1,\dots,q_s).
\]
Then:
\begin{enumerate}[label=(\arabic*)]
\item the image of $\rho_{m;S}$ is the compatibility locus
\[
\left\{(a_1,\dots,a_s)\in\prod_i \ZZ/(q_i\ZZ):
a_i\equiv a_j\pmod{\gcd(q_i,q_j)}\text{ for all }i,j\right\};
\]
\item the image has cardinality $Q$;
\item every fiber has cardinality $M/Q$.
\end{enumerate}
\end{theorem}

\begin{proof}
Under Proposition~\ref{thm:finite-resolution-mod}, the map $\rho_{m;S}$
identifies with
\[
\ZZ/(M\ZZ)\longrightarrow \prod_{i=1}^s \ZZ/(q_i\ZZ),
\qquad
\overline{t}\longmapsto ([t]_{q_1},\dots,[t]_{q_s}).
\]
Its image is the usual compatibility locus from the noncoprime Chinese
remainder theorem, proving (1). Two residue classes $\overline{t}$ and
$\overline{u}$ have the same image if and only if $q_i\mid (t-u)$ for every
$i$, equivalently $Q\mid (t-u)$. Thus the kernel has size $M/Q$, which proves
(3). Since the domain has size $M$, the image has size $Q$.
\end{proof}

\begin{corollary}
\label{cor:joint-readout-pairwise-coprime}
If $q_1,\dots,q_s$ are pairwise coprime, then $\rho_{m;S}$ is surjective and
every fiber has cardinality
\[
\frac{F_{m+2}}{q_1\cdots q_s}.
\]
\end{corollary}

\begin{proof}
If the $q_i$ are pairwise coprime, then $Q=q_1\cdots q_s$ and the compatibility
conditions in Theorem~\ref{thm:finite-signature-joint-readout-image-fiber}(1)
are vacuous.
\end{proof}

\begin{definition}
\label{def:proper-shadow-complexity}
A \emph{proper shadow family} at depth $m$ is a finite family of proper
divisors $q_1,\dots,q_s$ of $F_{m+2}$. Define $\sigma(m)$ to be the least $s$
for which the corresponding joint readout is injective. If no such family
exists, set $\sigma(m)=\infty$.
\end{definition}

\begin{theorem}
\label{thm:proper-shadow-compressibility}
Let
\[
F_{m+2}=\prod_{i=1}^r p_i^{a_i}.
\]
Then:
\begin{enumerate}[label=(\arabic*)]
\item a proper shadow family $\{q_1,\dots,q_s\}$ is injective if and only if
$\lcmop(q_1,\dots,q_s)=F_{m+2}$;
\item $\sigma(m)=\infty$ if $r=1$;
\item $\sigma(m)=2$ if $r\ge 2$.
\end{enumerate}
\end{theorem}

\begin{proof}
Part (1) follows from
Theorem~\ref{thm:finite-signature-joint-readout-image-fiber}(3): the joint
readout is injective exactly when its fiber size is $1$, that is, when the
least common multiple of the moduli equals $F_{m+2}$.

If $r=1$, then $F_{m+2}$ is a prime power and every proper divisor misses the
full exponent, so the least common multiple of proper divisors is always still
proper. By part (1), no proper shadow family is injective.

If $r\ge 2$, no single proper divisor can be injective. On the other hand,
\[
q:=\frac{F_{m+2}}{p_1^{a_1}},
\qquad
q':=\frac{F_{m+2}}{p_2^{a_2}}
\]
are proper divisors with $\lcmop(q,q')=F_{m+2}$, so part (1) gives
$\sigma(m)=2$.
\end{proof}

\begin{definition}
\label{def:prime-power-shadow-complexity}
A proper shadow family at depth $m$ is \emph{prime-power} if every modulus is a
proper prime-power divisor of $F_{m+2}$. Define $\sigma_{\mathrm{pp}}(m)$ to be
the least size of an injective prime-power shadow family. If no such family
exists, set $\sigma_{\mathrm{pp}}(m)=\infty$.
\end{definition}

\begin{theorem}
\label{thm:prime-power-shadow-complexity}
Write
\[
F_{m+2}=\prod_{i=1}^r p_i^{a_i}.
\]
Then
\[
\sigma_{\mathrm{pp}}(m)=
\begin{cases}
\infty, & r=1,\\
r, & r\ge 2.
\end{cases}
\]
In particular, whenever $X_m$ is CRT-split,
\[
\sigma_{\mathrm{pp}}(m)=\omega(F_{m+2}).
\]
\end{theorem}

\begin{proof}
Assume first that an injective prime-power shadow family exists. By
Theorem~\ref{thm:proper-shadow-compressibility}(1), its least common multiple
must equal $F_{m+2}$. Since every allowed modulus is a prime power divisor of
$F_{m+2}$, the family must contain the full factor $p_i^{a_i}$ for each $i$.
Thus every injective prime-power family has size at least $r$.

If $r=1$, the only prime power carrying the full exponent is $F_{m+2}$ itself,
which is not proper, so no such family exists. If $r\ge 2$, the family
$\{p_1^{a_1},\dots,p_r^{a_r}\}$ consists of proper divisors and has least
common multiple $F_{m+2}$, hence it is injective by
Theorem~\ref{thm:proper-shadow-compressibility}(1).
\end{proof}

\begin{corollary}
\label{cor:prime-power-shadow-divisor-lower-bound}
Let $n:=m+2$. Then
\[
\sigma_{\mathrm{pp}}(m)\ge \tau(n)-4.
\]
In particular, prime-power shadow complexity is unbounded along the Fibonacci
modulus chain.
\end{corollary}

\begin{proof}
If $\sigma_{\mathrm{pp}}(m)=\infty$, there is nothing to prove. Otherwise, by
Theorem~\ref{thm:prime-power-shadow-complexity},
$\sigma_{\mathrm{pp}}(m)=\omega(F_n)$. For each divisor
$d\mid n$ with $d\notin\{1,2,6,12\}$, Theorem
~\ref{thm:primitive-field-quotients} gives a prime $p_d$ with $\alpha(p_d)=d$.
Since $d\mid n$, Proposition~\ref{prop:rank-of-apparition-basic}(1) implies
$p_d\mid F_n$. Distinct divisors $d$ give distinct primes $p_d$, because the
rank of apparition of a prime is unique. Hence
$\omega(F_n)\ge \tau(n)-4$.
\end{proof}

\appendix

\section{First-birth tables}

The tables below are included only as arithmetic summaries of the main
theorems. They are direct consequences of
Theorems~\ref{thm:birth-spectrum-mobius}
and~\ref{thm:prime-power-apparition-ladders}; no additional computation enters
the proofs.

\begin{table}[ht]
\centering
\small
\begin{tabular}{rrrrr}
\toprule
$n$ & $F_n$ & $A(n)$ & $A_{\mathrm{pr}}(n)$ & $A_{\mathrm{pp}}(n)$ \\
\midrule
$2$  & $1$   & $0$  & $0$ & $0$ \\
$3$  & $2$   & $1$  & $1$ & $1$ \\
$4$  & $3$   & $1$  & $1$ & $1$ \\
$5$  & $5$   & $1$  & $1$ & $1$ \\
$6$  & $8$   & $2$  & $0$ & $2$ \\
$7$  & $13$  & $1$  & $1$ & $1$ \\
$8$  & $21$  & $2$  & $1$ & $1$ \\
$9$  & $34$  & $2$  & $1$ & $1$ \\
$10$ & $55$  & $2$  & $1$ & $1$ \\
$11$ & $89$  & $1$  & $1$ & $1$ \\
$12$ & $144$ & $10$ & $0$ & $2$ \\
\bottomrule
\end{tabular}
\caption{Initial values of the first-birth spectra. The exceptional indices $6$
and $12$ are prime-field empty but not prime-power empty.}
\label{tab:birth-spectra-small}
\end{table}

\begin{table}[ht]
\centering
\small
\begin{tabular}{rrrrr}
\toprule
$p$ & $\alpha(p)$ & $\nu_p(F_{\alpha(p)})$ & $\alpha(p^2)$ & $\alpha(p^3)$ \\
\midrule
$2$  & $3$  & $1$ & $6$   & $6$ \\
$3$  & $4$  & $1$ & $12$  & $36$ \\
$5$  & $5$  & $1$ & $25$  & $125$ \\
$7$  & $8$  & $1$ & $56$  & $392$ \\
$11$ & $10$ & $1$ & $110$ & $1210$ \\
$13$ & $7$  & $1$ & $91$  & $1183$ \\
$17$ & $9$  & $1$ & $153$ & $2601$ \\
\bottomrule
\end{tabular}
\caption{Initial prime-power ladder data. For the odd primes listed here one
has $\nu_p(F_{\alpha(p)})=1$, so the ladder begins with a single primitive
field layer before geometric growth.}
\label{tab:prime-power-ladders-small}
\end{table}

\section*{Declarations}

% Add any submission-specific Springer statements that apply here, including
% funding, author contributions, code availability, and, if needed, the
% required LLM disclosure when use exceeded copy editing.

\noindent\textbf{Funding.} No specific funding was received for this work.

\smallskip

\noindent\textbf{Conflict of interest.} The authors declare that they have no
conflict of interest.

\smallskip
\noindent\textbf{Data availability.} Data sharing not applicable to this
article as no datasets were generated or analyzed during the current study.

\bibliographystyle{unsrt}
\bibliography{../../2026_golden_ratio_driven_scan_projection_generation_recursive_emergence/references,../../2026_golden_ratio_driven_scan_projection_generation_recursive_emergence/references_fibonacci,../../2026_golden_ratio_driven_scan_projection_generation_recursive_emergence/references_algebra}

\end{document}
